\documentclass[10pt,a4paper]{amsart}
\usepackage{amsmath,amssymb,amsthm,graphicx}
\usepackage{verbatim}
\newenvironment{reference}{\comment}{\endcomment}
%\newenvironment{reference}{}{}
\newenvironment{slogan}{\comment}{\endcomment}
\newenvironment{history}{\comment}{\endcomment}
% For commutative diagrams we use Xy-pic
\usepackage[all]{xy}

% We use 2cell for 2-commutative diagrams.
\xyoption{2cell}
% Theorem environments.
%
\theoremstyle{plain}
\newtheorem{theorem}[subsection]{Theorem}
\newtheorem{proposition}[subsection]{Proposition}
\newtheorem{lemma}[subsection]{Lemma}

\theoremstyle{definition}
\newtheorem{definition}[subsection]{Definition}
\newtheorem{example}[subsection]{Example}
\newtheorem{exercise}[subsection]{Exercise}
\newtheorem{situation}[subsection]{Situation}

\theoremstyle{remark}
\newtheorem{remark}[subsection]{Remark}
\newtheorem{remarks}[subsection]{Remarks}

\numberwithin{equation}{subsection}

% Macros
%
\def\lim{\mathop{\mathrm{lim}}\nolimits}
\def\colim{\mathop{\mathrm{colim}}\nolimits}
\def\Spec{\mathop{\mathrm{Spec}}}
\def\Hom{\mathop{\mathrm{Hom}}\nolimits}
\def\Ext{\mathop{\mathrm{Ext}}\nolimits}
\def\SheafHom{\mathop{\mathcal{H}\!\mathit{om}}\nolimits}
\def\SheafExt{\mathop{\mathcal{E}\!\mathit{xt}}\nolimits}
\def\Sch{\mathit{Sch}}
\def\Mor{\mathop{\mathrm{Mor}}\nolimits}
\def\Ob{\mathop{\mathrm{Ob}}\nolimits}
\def\Sh{\mathop{\mathit{Sh}}\nolimits}
\def\NL{\mathop{N\!L}\nolimits}
\def\CH{\mathop{\mathrm{CH}}\nolimits}
\def\proetale{{pro\text{-}\acute{e}tale}}
\def\etale{{\acute{e}tale}}
\def\QCoh{\mathit{QCoh}}
\def\Ker{\mathop{\mathrm{Ker}}}
\def\Im{\mathop{\mathrm{Im}}}
\def\Coker{\mathop{\mathrm{Coker}}}
\def\Coim{\mathop{\mathrm{Coim}}}

% Boxtimes
%
\DeclareMathSymbol{\boxtimes}{\mathbin}{AMSa}{"02}

%
% Macros for moduli stacks/spaces
%
\def\QCohstack{\mathcal{QC}\!\mathit{oh}}
\def\Cohstack{\mathcal{C}\!\mathit{oh}}
\def\Spacesstack{\mathcal{S}\!\mathit{paces}}
\def\Quotfunctor{\mathrm{Quot}}
\def\Hilbfunctor{\mathrm{Hilb}}
\def\Curvesstack{\mathcal{C}\!\mathit{urves}}
\def\Polarizedstack{\mathcal{P}\!\mathit{olarized}}
\def\Complexesstack{\mathcal{C}\!\mathit{omplexes}}
% \Pic is the operator that assigns to X its picard group, usage \Pic(X)
% \Picardstack_{X/B} denotes the Picard stack of X over B
% \Picardfunctor_{X/B} denotes the Picard functor of X over B
\def\Pic{\mathop{\mathrm{Pic}}\nolimits}
\def\Picardstack{\mathcal{P}\!\mathit{ic}}
\def\Picardfunctor{\mathrm{Pic}}
\def\Deformationcategory{\mathcal{D}\!\mathit{ef}}
\usepackage{cleveref}
\crefname{originallemma}{original lemma}{original lemmas}
\Crefname{originallemma}{Original lemma}{Original lemmas}
\makeatletter
\newlabel{lemma-universal-module}{{lemma-universal-module}{source}{}{sourcepointer-1}{}}
\newlabel{lemma-universal-module@cref}{{[originallemma][0][]lemma-universal-module}{[1][0][]source}{}{}{}}
\makeatother
\title{Differential operators of order k are represented by a sheaf of principal parts}
\author{The Stacks Project Authors}
\date{stacks-project:modules.tex:local-snapshot-20261012; GNU FDL 1.2 or later}
\begin{document}
\maketitle
\noindent\textit{Editorial note.} All delivered mathematical text is the author's own, copied from the pinned Stacks Project chapter source sources/stacks-modules/pinned-original-source.tex; no statement or proof step was rewritten, reordered or removed.\par\medskip
\noindent\textit{Editorial note.} Preamble provenance: the Stacks Project keeps one single shared preamble for every chapter; this card takes it directly from the staged sources/stacks-modules/preamble.tex as three verbatim spans (lines 8-12, 14-18 and 154-217), so no derived file was written and nothing was added or redefined. The wrapper class and the base packages are supplied by the assembler.\par\medskip
\noindent\textit{Editorial note.} Bibliography provenance: neither the counted statement nor its proof contains a citation (this lane rescans the whole span before writing the card), so the card carries no bibliography and no entry was reconstructed.\par\medskip
\noindent\textit{Editorial note.} Selection provenance: this card was opened under the file-based control contract (CONTROL.md), with the mechanical pre-card gate run on its statement and proof interval before the card was written.\par\medskip
\noindent\textit{Allowed standard prerequisites (editorial).}\begin{itemize}
\item The free module on the sets of local sections over a sheaf of rings and its shorthand notation, taken from the construction of the universal module of differentials. The notation is reproduced verbatim as a context span of this card (source lines 4984-4997); the fact that it comes from the proof of a specific lemma of this chapter is recorded in the navigation-pointer stub. Provides the ambient module on which the relations defining the sheaf of principal parts are imposed.
\item The notion of a differential operator of order k and the relations characterising it in terms of the operators obtained by composing with multiplications. The author's section on finite order differential operators, in which the counted statement occurs. Supplies the relations that the representing module must quotient out.
\end{itemize}\medskip
\ifdefined\hypertarget\hypertarget{sourcepointer-1}{}\fi
\noindent\textit{Source reference.} Original lemma lemma-universal-module: sources/stacks-modules/pinned-original-source.tex, source lines 5040--5040. Navigation pointer only: the counted text refers to the Stacks item lemma-universal-module, which belongs to another chapter of the project and is not part of this card. The pointer carries the target's own label as its number because this chapter's staged snapshot does not contain the other chapter's tag table; its content is not reproduced here.\par\medskip
\begin{lemma}
\label{lemma-module-principal-parts}
Let $X$ be a topological space.
Let $\mathcal{O}_1 \to \mathcal{O}_2$ be a map of sheaves of rings on $X$.
Let $\mathcal{F}$ be a sheaf of $\mathcal{O}_2$-modules.
Let $k \geq 0$. There exists a sheaf of $\mathcal{O}_2$-modules
$\mathcal{P}^k_{\mathcal{O}_2/\mathcal{O}_1}(\mathcal{F})$
and a canonical isomorphism
$$
\text{Diff}^k_{\mathcal{O}_2/\mathcal{O}_1}(\mathcal{F}, \mathcal{G}) =
\Hom_{\mathcal{O}_2}(
\mathcal{P}^k_{\mathcal{O}_2/\mathcal{O}_1}(\mathcal{F}), \mathcal{G})
$$
functorial in the $\mathcal{O}_2$-module $\mathcal{G}$.
\end{lemma}

During this proof, for any sheaf of sets $\mathcal{F}$ on $X$,
let us denote $\mathcal{O}_2[\mathcal{F}]$ the sheafification of the
presheaf $U \mapsto \mathcal{O}_2(U)[\mathcal{F}(U)]$ where this denotes
the free $\mathcal{O}_2(U)$-module on the set $\mathcal{F}(U)$.
For $s \in \mathcal{F}(U)$ we denote $[s]$ the corresponding section
of $\mathcal{O}_2[\mathcal{F}]$ over $U$. If $\mathcal{F}$ is a sheaf of
$\mathcal{O}_2$-modules, then there is a canonical map
$$
c : \mathcal{O}_2[\mathcal{F}] \longrightarrow \mathcal{F}
$$
which on the presheaf level is given by the rule
$\sum f_s[s] \mapsto \sum f_s s$. We will employ the short hand
$[s] \mapsto s$ to
describe this map and similarly for other maps below. Consider

\begin{proof}
The existence follows from general category theoretic arguments
(insert future reference here), but we will also give a direct
construction as this construction will be useful in the future proofs.
We will freely use the notation introduced in the proof of
Lemma \ref{lemma-universal-module}.
Given any differential operator $D : \mathcal{F} \to \mathcal{G}$
we obtain an $\mathcal{O}_2$-linear map
$L_D : \mathcal{O}_2[\mathcal{F}] \to \mathcal{G}$
sending $[m]$ to $D(m)$. If $D$ has order $0$
then $L_D$ annihilates the local sections
$$
[m + m'] - [m] - [m'],\quad
g_0[m] - [g_0m]
$$
where $g_0$ is a local section of $\mathcal{O}_2$ and $m, m'$
are local sections of $\mathcal{F}$. If $D$ has order $1$, then $L_D$
annihilates the local sections
$$
[m + m' - [m] - [m'],\quad
f[m] - [fm], \quad
g_0g_1[m] - g_0[g_1m] - g_1[g_0m] + [g_1g_0m]
$$
where $f$ is a local section of $\mathcal{O}_1$,
$g_0, g_1$ are local sections of $\mathcal{O}_2$, and
$m, m'$ are local sections of $\mathcal{F}$.
If $D$ has order $k$, then $L_D$ annihilates the local sections
$[m + m'] - [m] - [m']$, $f[m] - [fm]$, and the local sections
$$
g_0g_1\ldots g_k[m] - \sum g_0 \ldots \hat g_i \ldots g_k[g_im] + \ldots
+(-1)^{k + 1}[g_0\ldots g_km]
$$
Conversely, if $L : \mathcal{O}_2[\mathcal{F}] \to \mathcal{G}$ is an
$\mathcal{O}_2$-linear map annihilating all the local sections
listed in the previous sentence, then $m \mapsto L([m])$ is a
differential operator of order $k$. Thus we see that
$\mathcal{P}^k_{\mathcal{O}_2/\mathcal{O}_1}(\mathcal{F})$
is the quotient of $\mathcal{O}_2[\mathcal{F}]$
by the $\mathcal{O}_2$-submodule generated by these local sections.
\end{proof}

\end{document}
